\section{Time for a Second Renunciation}

Several of \textbf{Miguel Serrano}'s books had been made available in English years before anyone in the Anglosphere was aware of his ties to the extreme right wing. His reputation as an ambassador from Chile and his association with several iconic figures such as \textbf{Hermann Hesse}, \textbf{Carl Jung}, and the \textbf{Dalai Lama} made him influential in certain rather small circles. Serrrano was intelligent, sophisticated, knowledgeable in esoteric matters, qualities which rendered his interviews particularly interesting and incisive. In The Serpent of Paradise, which describes his search for ancient knowledge in India.

One of his visits was with \textbf{Jiddu Krishnamurti} who, as a youth, was proclaimed the Messiah by certain elements of the Theosophical Society. Krishnamurti rejected that role and all the benefits that would have accrued to him. Serrano writes:

\begin{quotex}
Krishnamurti has been an extraordinarily valuable example to the world because he resisted the supreme temptation of power and wealth, and he renounced all that might have been due him as a Messiah. He continually resisted the adoration of psychopaths, most of whom were rich and idle, the sort of people who go about the world, anxious to throw themselves at the feet of the first semidivine figure they encounter.
\end{quotex}

This is because men are quite reluctant to make the effort to change themselves and rest content in sitting under the aura of some other figure. The wealthy, in particular, are accustomed to pay others for their menial work while they get tax deductions for their lavish gifts. But even the not so wealthy are so prone, yet rest content with books and tapes.

I used to attend one such study group at a public park in Ft Lauderdale in the 1980s. One evening, the group leader was excited by an item in the newspaper that day that claimed there was a sufficient supply of food in the world for all mankind, but that the proper distribution was lacking. I pointed out that this was ``bad'' news. Solving a problem with technology is one thing, but modifying human nature is vastly different. But the leader was full of ``love'' for mankind and assumed everyone would share their food. That was my last meeting, as I was no longer welcome there.

Serrano speculates that Krishnamurti was ``naturally made for love, for surrender, and for obedience,'' that is, he was essentially a Bhakti Yogi. He then forced himself into a Jnana Yogi, and ``intellectual path, which is hard and merciless.'' Krishnamurti's essential message is this:

\begin{quotex}
There is no Master, there are no sacred books and there is no tradition. Nobody can teach anybody anything. Nobody should listen to anyone else; there is no sense in following.
\end{quotex}


This we can relate to the First Trial\footnote{\url{https://gornahoor.net/?p=1902}}, purging one's mind of all opinions. Instead, Krishnamurti recommends attention: learning how to look at the world or to listen to one's son, wife, or friend. Full attention leads to the second Trial which is reflected in his attitude to death.

\begin{quotex}
Whoever lives in the present cannot fear death, because he has put himself totally in the act of living. Thus, even when he dies, he will have no fear of death, since in the very act of dying he will have gained totality, perhaps for the first time in his life. This dying person will respond totally to the request; he will give himself entirely, with all his life to death. To die like this is an act of love.
\end{quotex}


This lead to the Third Trial. Krishnamurti says:

\begin{quotex}
To love, or even to kill, is to put oneself wholly into an action and to gain an eternity in the present tense. The notion of fear arises because we are not total: a bit of us lives in the past and a part of us in the future.
\end{quotex}


This prompts Serrano to pose this question: ``\emph{Is the act of murder, or the commission of some crime, as pure an act as that of loving?}'' Krishnamurti answers:

\begin{quotex}
Yes, but only insofar as that action leaves no trace on the mind, only insofar as the mind remains untouched by it. For all actions should take place in that way. Indeed, love should leave no traces once it has been lived, or, like a crime, been committed.
\end{quotex}


However, it is unlikely that Krishnamurti had ever loved or killed, a paradox noted by Serrano. While renouncing all prior teachings and convention, Krishnamurti nevertheless lived quite conventionally the life of a Brahman from the South of India. Krishnamurti renounced his role as a Messiah. Serrano speculates that a second renunciation is necessary.

\begin{quotex}
To be able to advance and to flourish, like the rose in his room ... he will have to love or to kill... He will have to become a total man and to descend to the ways of man. \emph{In short, it is time for a second renunciation}.
\end{quotex}



\flrightit{Posted on 2011-06-05 by Cologero}

\begin{center}* * *\end{center}

\begin{footnotesize}
\begin{sffamily}

\texttt{logres on 2011-08-14 at 19:07 said:}

Reading about Krishmanurti -- this post is fascinating. What does a ``second renunciation'' refer to? The renouncing of Messiahood? Or a willingness to engage in the affairs of men? Or is it both?

\hfill

\texttt{Cologero on 2011-08-14 at 23:41 said:}

Perhaps I should have included the entire paragraph:

\begin{quotex}\footnotesize
To be able to advance and to flourish, like the rose in his room, he will have to renounce his lectures, or he will have to love or to kill, or he will have to stop being the Messiah which in spite of himself he still is. He will have to become a total man and to descend to the ways of man. In short, it is time for a second renunciation.
\end{quotex}


\hfill

\texttt{Balder on 2022-08-23 at 12:18 said:}

About the renunciation of all forms and traditions etc., I think it is possible and useful only for a very advanced initiate-esoterist; it is inadvisable to recommend the pathless path taken by Khrisnamurti or the like for any and all without reserve, since the rejection of forms and traditions will easily lead the ``passive flock'' and multitudes into dissolution and confusion.

Although the Formless Source is really the one thing needful, one cannot straight forwards climb into Heaven without ladders, and this is why traditions, rituals/sacraments, and forms have their own useful place, if we do not fall into a trap to consider they have meaning ex cathedra without the imposition of human mediator-like intention into the acts of formal manifestation, invocation, evocation etc.

\hfill

\texttt{Balder on 2022-08-23 at 12:32 said:}

Oh, and I forgot to mention that the best apology for tradition in general is the fact that it has developed over the millenia which has made it more resilient towards individual error; yet, this is unilaterally a right hand path method, the impersonal impulse where the person is not allowed to shine its unique light.

This on the other brings to surface the importance of a left hand path-right hand path synthesis, where the ``lower'' individuality is also appreciated and one does not strive towards total impersonalisation. There is no need ``to get rid of the ego'', since once it is purified of lower impulses, it is simply a very useful and neutral instrument for the working out of actions in the threefold universe.

Regarding Khrisnamurti it is clear that he followed a very individual, even anarchistic path in regards tradition, and this alone in a sense makes him a LHP figure Yet as mentioned, he lived outwardly a life of a RHP Brahmin.

\end{sffamily}
\end{footnotesize}

